\documentclass[UTF8,a4paper,11pt]{ctexart}
\usepackage[a4paper,margin=2.3cm]{geometry}
\usepackage{amsmath,amssymb,mathtools}
\usepackage{enumitem}
\usepackage[hidelinks]{hyperref}
\usepackage{fancyhdr}
\setlist[itemize]{leftmargin=2.2em,itemsep=0.25em,topsep=0.35em}
\setlength{\parindent}{2em}\setlength{\parskip}{0.35em}
\pagestyle{fancy}\fancyhf{}\fancyhead[L]{离散数学与结构（I）}\fancyhead[R]{Lec 11}\fancyfoot[C]{\thepage}\setlength{\headheight}{14pt}
\title{分裂域、有限域与有限域上的多项式因式分解}\author{zhangjunshuo}\date{2026 年 10 月}
\begin{document}
\maketitle
\begin{quote}
本讲根据 2026 年 10 月 20 日的两段课堂转写、三页手写原稿和课程教材整理。本讲紧接 Lec10 末尾的域扩张与分裂域，先用域扩张语言（次数、三种造域方式、代数元与极小多项式）打底，用塔律证明倍立方与三等分角不可作；再证明分裂域存在且在同构意义下唯一；然后用它证明有限域 $\mathbb F_{p^d}$ 的存在唯一性，把 $x^{p^d}-x$ 分解成不可约多项式之积、数出每个次数的不可约多项式个数，并给出有限域上多项式因式分解的三步算法（去平方因子、按次数分组、Cantor--Zassenhaus 同次数拆分）。

每个分节开头都先交代本节用到的记号；所有符号在首次使用处给出定义，不依赖读者记得前面讲次的写法。
\end{quote}
\tableofcontents\clearpage

\section{域扩张回顾：次数、三种造域方式与代数元}


Lec10 的结尾用理想构造了域，并引出域扩张与分裂域。这一节先把那批工具整理成 Lec11 反复使用的语言：\textbf{域扩张的次数}、\textbf{三种造域方式}、\textbf{域同态为什么总是单射}，以及\textbf{代数元与极小多项式}。分裂域、有限域、因式分解三部分都建立在这些语言上，所以这一节的控制目标是：后面出现 $[K:F]$、$F(\alpha)$、极小多项式时，不再需要回头查定义。

\textbf{本节用到的记号：} $F,K,E,H$ 为域，$F\subseteq K$ 表示 $F$ 是 $K$ 的子域；$K/F$ 表示扩域（读作“$K$ over $F$”）；$[K:F]=\dim_F K$ 为扩张次数；$F[x]$ 为多项式环，$(f)$ 为 $f$ 生成的主理想；$F(\alpha)$ 为 $K$ 中包含 $F$ 与 $\alpha$ 的最小子域；$\operatorname{Frac}(R)$ 为分式域；$\mathbb F_p=\mathbb Z/p\mathbb Z$；$\operatorname{char}(F)$ 为特征。

\subsection{域扩张与次数}


设 $F\subseteq K$ 都是域，且 $F$ 的加法、乘法与 $K$ 一致。此时 $K$ 自然地是 $F$ 上的\textbf{线性空间}：向量加法就是 $K$ 的加法，标量乘法就是 $F$ 中元素乘 $K$ 中元素。它的维数记作

\[
[K:F]=\dim_F K,
\]

称为扩张 $K/F$ 的\textbf{次数}。当 $[K:F]<\infty$ 时称 $K/F$ 为\textbf{有限扩张}，否则称为\textbf{无限扩张}。

这个“把扩域看成线性空间”的视角是整讲的计算引擎：次数可乘（塔律）、元素有极小多项式、有限域大小是 $p^d$，全部从它推出来。

\subsection{三种造域方式（回顾）}


Lec10 已经给出三条从已知域造出新域的路线，这里只列结论并把课上强调的细节补全。

\textbf{路线一：商掉一个不可约多项式。} 设 $F$ 是域、$f\in F[x]$ 不可约、$\deg f=n\ge1$，则 $(f)$ 是极大理想，商环

\[
F[x]/(f)
\]

是域。它的每个元素有唯一代表元

\[
a_0+a_1\bar x+\cdots+a_{n-1}\bar x^{\,n-1},\qquad a_i\in F,
\]

其中 $\bar x$ 是 $x$ 的像。因此 $[F[x]/(f):F]=n$，且 $|F[x]/(f)|=|F|^n$。

\textbf{路线二：取分式域。} 从整环 $R$ 出发，把除法补进去得到 $\operatorname{Frac}(R)$。本讲最常用的是 $R=F[x]$，此时

\[
\operatorname{Frac}(F[x])=F(x)=\left\{\frac{g(x)}{h(x)}:g,h\in F[x],\ h\ne0\right\},
\]

称为 $F$ 上的\textbf{有理函数域}。它是 $F$ 的无限扩张。

\textbf{路线三：在一个更大的域里添一个元素。} 设 $K/F$ 是扩域、$\alpha\in K$。记 $F(\alpha)$ 为 $K$ 中包含 $F$ 与 $\alpha$ 的\textbf{最小子域}。它是 $F$ 的扩域，且 $[F(\alpha):F]$ 可以是有限，也可以是无限。

三条路线的关系：路线一造出“形式上的根”，路线三在具体扩域里“取一个真实的根”，路线二则负责“添分母”。下一小节说明路线一与路线三在“根”这件事上完全一致。

\subsection{添一个根：$F[x]/(f)\cong F(\alpha)$}


\textbf{定理。} 设 $F$ 是域、$f\in F[x]$ 不可约、$K/F$ 是扩域、$\alpha\in K$ 满足 $f(\alpha)=0$。则存在域同构

\[
F[x]/(f)\ \cong\ F(\alpha),
\]

它保持 $F$ 中元素不动，并把 $\bar x$ 送到 $\alpha$。

\textbf{证明。} 定义求值映射

\[
\varphi:F[x]\longrightarrow F(\alpha),\qquad g(x)\longmapsto g(\alpha).
\]

它是环同态：$\varphi(g+h)=g(\alpha)+h(\alpha)$、$\varphi(gh)=g(\alpha)h(\alpha)$，且 $\varphi$ 在 $F$ 上是恒等映射。下面三步。

\textbf{(1) 核是 $(f)$。} 显然 $f\in\ker\varphi$，所以 $(f)\subseteq\ker\varphi$。设 $g\in\ker\varphi$，即 $g(\alpha)=0$。因为 $f$ 不可约，$(f)$ 是极大理想，而 $\ker\varphi$ 是包含 $(f)$ 的理想，所以 $\ker\varphi$ 只能是 $(f)$ 或整个 $F[x]$。后者不可能，因为 $\varphi(1)=1\ne0$ 说明 $1\notin\ker\varphi$。故 $\ker\varphi=(f)$。

（同一件事也可以这样看：取 $d=\gcd(f,g)$。因为 $d$ 是 $f,g$ 的 $F[x]$-线性组合，所以 $d(\alpha)=0$；但 $f$ 不可约说明 $d$ 是常数或 $f$ 的相伴元，而常数不可能有 $d(\alpha)=0$，故 $f\mid g$，即 $g\in(f)$。这种写法在 Lec10 用过，等价于“极大理想”论证。）

\textbf{(2) 满射。} 任取 $\beta\in F(\alpha)$。由 $F(\alpha)$ 的定义，它是 $F$ 与 $\alpha$ 生成的域，其中每个元素都可以写成 $\alpha$ 的有理式；又因 $f(\alpha)=0$，用带余除法可以把 $\alpha$ 的高次幂降到低于 $n=\deg f$，所以存在多项式 $g\in F[x]$ 使 $\beta=g(\alpha)$。故 $\varphi$ 满。

\textbf{(3) 由第一同构定理。} $\varphi$ 诱导 $F[x]/\ker\varphi\cong\operatorname{im}\varphi$，即 $F[x]/(f)\cong F(\alpha)$；$\bar x$ 的像为 $\varphi(x)=\alpha$。$\square$

\textbf{这条定理在说什么。} 它把两件事绑定在一起：\textbf{“在抽象商环里添一个形式根”与“在具体扩域里取一个真实根”得到同一个域}。所以以后想造 $n$ 次扩域，只要找一个 $n$ 次不可约多项式即可。

\textbf{一个容易忽略的前提。} $f$ 的不可约性不能省。若 $f$ 可约，$(f)$ 不是极大理想，$F[x]/(f)$ 只是环而不是域。例如 $\mathbb R[x]/(x^2-1)$ 有零因子，它不是域。

\textbf{同构不唯一。} 若 $\alpha,\alpha'$ 是同一个不可约多项式 $f$ 的两个不同根，则 $F(\alpha)\cong F[x]/(f)\cong F(\alpha')$，但把 $\alpha$ 送到 $\alpha'$ 还是送到别的根，得到的是不同的同构。这正是分裂域“同构意义下唯一”的伏笔。

\subsection{域同态只有两种：单射或零映射}


\textbf{命题。} 设 $\varphi:F\to F'$ 是域同态，则 $\varphi$ 要么是单射，要么是零映射（把所有元素送到 $0$）。

\textbf{证明。} $\ker\varphi$ 是 $F$ 的理想。域只有两个理想：$(0)$ 和 $F$ 本身。若 $\ker\varphi=(0)$，则 $\varphi$ 是单射；若 $\ker\varphi=F$，则 $\varphi$ 是零映射。$\square$

这条命题的用法：\textbf{两个域之间的非零同态一定是嵌入（单射）}。因此本讲后面说“保持 $F$ 不动的同态”时，自动就是“把 $F$ 的扩域嵌入另一个扩域”。

\textbf{推论。} 若 $\varphi:F\to F'$ 是非零域同态，则 $\varphi(F)$ 是与 $F$ 同构的子域。

\subsection{代数元、超越元与极小多项式}


设 $K/F$ 是扩域、$\alpha\in K$。

\begin{itemize}
\item 若存在\textbf{非零}多项式 $g\in F[x]$ 使 $g(\alpha)=0$，称 $\alpha$ 在 $F$ 上\textbf{代数}；
\item 否则称 $\alpha$ 在 $F$ 上\textbf{超越}。
\end{itemize}


\textbf{等价刻画。} $\alpha$ 在 $F$ 上代数 $\iff$ $[F(\alpha):F]<\infty$。

\textbf{证明思路。} 若 $\alpha$ 代数，取一个非零多项式 $g$ 使 $g(\alpha)=0$，设 $\deg g=n$。则 $1,\alpha,\dots,\alpha^n$ 在 $F$ 上线性相关（因为 $g$ 给出一个非平凡线性关系），于是 $[F(\alpha):F]\le n<\infty$（更精细地，次数等于极小多项式的次数，见下）。反过来，若 $[F(\alpha):F]=n<\infty$，则 $1,\alpha,\dots,\alpha^n$ 这 $n+1$ 个元素在 $n$ 维空间里线性相关，给出非零多项式 $g$ 使 $g(\alpha)=0$。$\square$

\textbf{极小多项式。} 设 $\alpha$ 在 $F$ 上代数。在所有满足 $g(\alpha)=0$ 的非零多项式中，取次数最小者，再除以首项系数使它\textbf{首一}，得到的多项式称为 $\alpha$ 在 $F$ 上的\textbf{极小多项式}，记作 $m_{\alpha,F}(x)$（或简记 $m_\alpha$）。它有两个基本性质：

\textbf{(1) 极小多项式不可约。} 若 $m_\alpha=gh$ 且 $g,h$ 次数都小于 $\deg m_\alpha$，则由 $0=m_\alpha(\alpha)=g(\alpha)h(\alpha)$ 及 $F(\alpha)$ 是域（无零因子）得 $g(\alpha)=0$ 或 $h(\alpha)=0$，与极小性矛盾。故 $m_\alpha$ 不可约。

\textbf{(2) 极小多项式整除任何以 $\alpha$ 为根的多项式。} 设 $g(\alpha)=0$。作带余除法 $g=q\,m_\alpha+r$，$\deg r<\deg m_\alpha$。代入 $\alpha$ 得 $r(\alpha)=0$，由极小性 $r=0$，故 $m_\alpha\mid g$。

\textbf{推论（次数公式）。} $[F(\alpha):F]=\deg m_{\alpha,F}$。

\textbf{证明。} 由上面的定理取 $f=m_\alpha$（它不可约、以 $\alpha$ 为根），得 $F(\alpha)\cong F[x]/(m_\alpha)$，而后者在 $F$ 上的维数是 $\deg m_\alpha$。$\square$

\textbf{超越元的情形。} 若 $\alpha$ 超越，则映射

\[
F[x]\longrightarrow F(\alpha),\qquad g(x)\longmapsto g(\alpha)
\]

没有非零核元素，所以是单射；它自然延拓到分式域得到

\[
\operatorname{Frac}(F[x])=F(x)\ \cong\ F(\alpha),\qquad \frac{g(x)}{h(x)}\longmapsto\frac{g(\alpha)}{h(\alpha)}.
\]

（良定义：$h\ne0$ 且 $\alpha$ 超越，故 $h(\alpha)\ne0$；与分数表示无关由交叉相乘验证；单射由 $g(\alpha)=0\Rightarrow g=0$；满射由 $F(\alpha)$ 中元素都是 $\alpha$ 的有理式。）于是 $[F(\alpha):F]=\infty$。这正是“超越 $=$ 无限次扩张”的含义。

\textbf{有限扩张必代数。} 若 $[K:F]<\infty$，则任意 $\alpha\in K$ 都满足 $[F(\alpha):F]\le[K:F]<\infty$，故 $\alpha$ 在 $F$ 上代数。\textbf{反之不然}：存在无限次代数扩张（例如 $\mathbb Q$ 在 $\mathbb C$ 中的代数闭包），但本讲不展开。

\subsection{几个必须记住的例子}


\begin{itemize}
\item $\alpha=\sqrt2$ 在 $\mathbb Q$ 上代数，$m_\alpha=x^2-2$，$[\mathbb Q(\sqrt2):\mathbb Q]=2$。
\item $\alpha=\sqrt[3]{2}$ 在 $\mathbb Q$ 上代数，$m_\alpha=x^3-2$（在 $\mathbb Q$ 上不可约：有理根只能是 $\pm1,\pm2$，代入都不为 $0$），$[\mathbb Q(\sqrt[3]{2}):\mathbb Q]=3$。
\item $\pi$ 在 $\mathbb Q$ 上超越（Lindemann 1882），$[\mathbb Q(\pi):\mathbb Q]=\infty$。本课程不给证明，只作为“无限扩张”的例子使用。
\item 设 $F$ 是域、$x$ 是形式变量，则 $x$ 在 $F$ 上超越，$F(x)$ 是 $F$ 的无限扩张（路线二）。
\item $\mathbb C/\mathbb R$：$[\mathbb C:\mathbb R]=2$，$m_i=x^2+1$。
\end{itemize}


\subsection{本节的勘误与提醒}


\begin{itemize}
\item \textbf{$F[x]/(f)\cong F(\alpha)$ 不是“等于”。} 商环的元素是陪集，$F(\alpha)$ 的元素是扩域中的元素，二者之间只有同构；同构由 $\bar x\mapsto\alpha$ 决定。
\item \textbf{$f$ 不可约是商环成为域的前提；$f(\alpha)=0$ 是把 $\bar x$ 送到 $\alpha$ 的前提。} 两个条件缺一不可。
\item \textbf{极小多项式必须首一}，否则“唯一”不成立（$2x^2-4$ 与 $x^2-2$ 只差常数倍）。
\item \textbf{极小多项式的次数等于扩张次数，不是“不超过”。} 上文的线性相关只给出 $\le n$，真正等号来自 $F(\alpha)\cong F[x]/(m_\alpha)$。
\item \textbf{“代数”与“有限扩张”不是同一个概念。} 代数元给出有限扩张，但扩域可以是无限次代数扩张；有限扩张则一定是代数扩张。
\item 课堂转写里“$K$ 是一个 $F$-linear space”“$[K:F]$ 是维数”这类句子识别得很乱，以本节定义为准。
\end{itemize}

\section{塔律与尺规作图三不可能}


上一节把域扩张的次数定义成线性空间维数。这一节先给出次数的乘法公式（\textbf{塔律}），它让次数变得可计算；然后把“尺规作图”翻译成域扩张，用它证明古希腊三大作图难题中的两个——\textbf{倍立方}与\textbf{三等分角}——在尺规作图下不可能。第三个难题\textbf{化圆为方}需要 $\pi$ 的超越性，本讲只提结论。

\textbf{本节用到的记号：} $F\subseteq K\subseteq H$ 为域塔；$[H:F]$ 为扩张次数；$F(\alpha)$ 为添一个元素得到的最小子域；$m_{\alpha,F}$ 为极小多项式；$\mathbb Q$ 为有理数域；尺规作图的可作点、可作长度见下文定义。

\subsection{次数乘法公式（塔律）}


\textbf{定理（塔律）。} 设 $F\subseteq K\subseteq H$ 是域塔，则

\[
[H:F]=[H:K]\,[K:F].
\]

当两边出现无限时，按“$\infty\cdot$ 正数 $=\infty$”理解；有限情形的证明如下。

\textbf{证明（有限情形）。} 设 $[K:F]=m$，取 $K$ 在 $F$ 上的一组基 $e_1,\dots,e_m$；设 $[H:K]=n$，取 $H$ 在 $K$ 上的一组基 $h_1,\dots,h_n$。考虑 $mn$ 个元素

\[
\{\,e_i h_j:1\le i\le m,\ 1\le j\le n\,\}\subseteq H.
\]

\textbf{第一步：它们张成 $H$（作为 $F$-线性空间）。} 任取 $u\in H$。因为 $\{h_j\}$ 是 $H$ 在 $K$ 上的基，存在 $b_j\in K$ 使

\[
u=\sum_{j=1}^{n} b_j h_j .
\]

又因为 $\{e_i\}$ 是 $K$ 在 $F$ 上的基，每个 $b_j$ 可写成 $b_j=\sum_{i=1}^{m}a_{ij}e_i$，$a_{ij}\in F$。代入得

\[
u=\sum_{j=1}^{n}\Bigl(\sum_{i=1}^{m}a_{ij}e_i\Bigr)h_j=\sum_{i=1}^{m}\sum_{j=1}^{n}a_{ij}\,(e_i h_j),
\]

所以 $\{e_ih_j\}$ 张成 $H$。

\textbf{第二步：它们 $F$-线性无关。} 设 $\sum_{i,j}c_{ij}e_ih_j=0$，$c_{ij}\in F$。按 $j$ 合并：

\[
\sum_{j=1}^{n}\Bigl(\sum_{i=1}^{m}c_{ij}e_i\Bigr)h_j=0.
\]

括号里的元素属于 $K$，而 $\{h_j\}$ 在 $K$ 上线性无关，所以每个括号为 $0$，即 $\sum_i c_{ij}e_i=0$。再由 $\{e_i\}$ 在 $F$ 上线性无关得所有 $c_{ij}=0$。

于是 $\{e_ih_j\}$ 是 $H$ 在 $F$ 上的一组基，共有 $mn$ 个元素，故 $[H:F]=mn=[H:K][K:F]$。$\square$

\textbf{无限情形的说明。} 若 $[K:F]=\infty$，则 $K$ 中存在无限多个 $F$-线性无关的元素，它们在更大的 $H$ 中仍然 $F$-线性无关，故 $[H:F]=\infty$；若 $[H:K]=\infty$，同理 $H$ 中存在无限多个 $K$-线性无关的元素，它们对 $F$ 也线性无关，故 $[H:F]=\infty$。两种情形两边都是 $\infty$。

\textbf{推论一（整除性）。} 若 $F\subseteq K\subseteq H$ 且 $[H:F]<\infty$，则 $[K:F]$ 与 $[H:K]$ 都整除 $[H:F]$。

这条推论是这一节证明“不可能”的核心工具：\textbf{中间域的次数必须整除总次数。}

\textbf{推论二（链式法则）。} 若 $F\subseteq F(\alpha_1)\subseteq F(\alpha_1,\alpha_2)\subseteq\cdots$，则总次数等于每一步次数之积。

\subsection{尺规作图：把几何翻译成域}


取平面直角坐标系，给定两个点 $(0,0)$ 与 $(1,0)$。\textbf{尺规作图}允许的操作只有三种：

\begin{itemize}
\item 过两个已作点作直线；
\item 以一个已作点为圆心、以它到另一个已作点的距离为半径作圆；
\item 取两条已作直线、一条直线与一个圆、或两个圆的交点。
\end{itemize}


\textbf{可作点}指从 $(0,0),(1,0)$ 出发，经过有限次上述操作能得到的点；两点距离称为\textbf{可作长度}。

\textbf{关键观察（每一步最多二次）。} 设当前所有已作点的坐标都属于某个域 $F$。新点 $P$ 的产生只有三种情况：

\begin{itemize}
\item 直线与直线相交：交点坐标是 $F$ 上线性方程组的解，仍在 $F$ 中；
\item 直线与圆相交：把直线方程代入圆方程，得到一个\textbf{二次方程}，其系数在 $F$ 中，故交点坐标属于 $F$ 的至多二次扩张 $F(\sqrt{\Delta})$；
\item 两圆相交：两圆方程相减得到一条直线（根轴），于是化归为“直线与圆”的情形。
\end{itemize}


所以每作一个新点，坐标所在的域 $F$ 至多扩大成 $F$ 的\textbf{二次扩张}：新域次数是原来的 $1$ 倍或 $2$ 倍。由塔律，若经过 $m$ 次真正扩大次数的步骤，则最终坐标域的扩张次数为

\[
[K:\mathbb Q]=2^{m}.
\]

\textbf{结论。} 若数 $x$ 是可作长度（特别地，若点 $(x,0)$ 可作），则 $x$ 属于某个满足 $[K:\mathbb Q]=2^m$ 的域 $K$，于是由推论一，

\[
[\mathbb Q(x):\mathbb Q]\ \big|\ 2^{m}.
\]

换句话说：\textbf{可作数的极小多项式次数必须是 $2$ 的幂。} 这是所有“不可能”证明的统一判据。

\subsection{倍立方不可能}


\textbf{问题。} 已知单位立方体（体积 $1$），要求用尺规作一个体积为 $2$ 的立方体，即作出长度 $\sqrt[3]{2}$。

\textbf{证明不可能。} 令 $\alpha=\sqrt[3]{2}$，它满足 $x^3-2=0$。多项式 $x^3-2$ 在 $\mathbb Q$ 上不可约：由有理根定理，它的有理根只能是 $\pm1,\pm2$，而 $1^3-2=-1$、$(-1)^3-2=-3$、$2^3-2=6$、$(-2)^3-2=-10$，都不为 $0$。因此

\[
[\mathbb Q(\sqrt[3]{2}):\mathbb Q]=\deg(x^3-2)=3.
\]

若 $\sqrt[3]{2}$ 可作，则 $3$ 必须整除 $2^m$——矛盾。$\square$

\subsection{三等分角不可能（以 $60^\circ$ 为例）}


\textbf{问题。} 给定任意角，要求用尺规把它三等分。我们只需证明\textbf{存在一个角不能三等分}：取 $60^\circ$。

$60^\circ$ 本身可作（作等边三角形即可）。若 $60^\circ$ 能三等分，则 $20^\circ$ 可作，从而 $c=\cos20^\circ$ 是可作长度。

\textbf{用三倍角公式翻译成多项式。} 对任意角 $\theta$，

\[
\cos 3\theta=4\cos^3\theta-3\cos\theta .
\]

取 $\theta=20^\circ$，则 $3\theta=60^\circ$、$\cos 3\theta=\tfrac12$，代入得

\[
4c^3-3c=\frac12,
\qquad\text{即}\qquad
8c^3-6c-1=0 .
\]

令 $u=2c$，则 $u^3-3u-1=0$。多项式 $x^3-3x-1$ 在 $\mathbb Q$ 上不可约：它是首一整系数三次多项式，有理根只能是 $\pm1$，而 $f(1)=1-3-1=-3$、$f(-1)=-1+3-1=1$，都不为 $0$。因此

\[
[\mathbb Q(u):\mathbb Q]=3,\qquad\text{从而}\quad [\mathbb Q(c):\mathbb Q]=3
\]

（因为 $\mathbb Q(u)=\mathbb Q(c)$，$u=2c$ 只是乘以有理数）。若 $c$ 可作，则 $3\mid 2^m$，矛盾。$\square$

\textbf{不要误读成“所有角都不能三等分”。} 例如 $90^\circ$ 可以三等分成 $30^\circ$；$60^\circ$ 不能。命题是“尺规作图不能三等分任意角”。

\subsection{化圆为方（只给结论）}


\textbf{问题。} 已知单位圆（面积 $\pi$），要求用尺规作一个面积等于 $\pi$ 的正方形，即作出长度 $\sqrt{\pi}$。

$\sqrt{\pi}$ 可作的一个必要条件是 $\sqrt{\pi}$ 在 $\mathbb Q$ 上代数，而 $\sqrt{\pi}$ 代数等价于 $\pi$ 代数（若 $\pi$ 是非零多项式 $p(x)$ 的根，则 $\sqrt\pi$ 是 $p(x^2)$ 的根；反之 $x^2-\pi$ 的根是 $\sqrt\pi$，若 $\sqrt\pi$ 代数则 $\pi=(\sqrt\pi)^2$ 代数）。但 Lindemann 在 1882 年证明了 \textbf{$\pi$ 是超越数}，所以 $\sqrt\pi$ 超越，不可能可作。本课程只引用这个结论，不证明。

\subsection{为什么“次数是 $2$ 的幂”只是必要条件}


上面的论证只说明：可作数必须满足 $[\mathbb Q(x):\mathbb Q]=2^k$。\textbf{反过来不成立}：一个次数为 $2^k$ 的数未必可作，因为尺规作图要求一座\textbf{每一步都是二次}的域塔，而不是一步到位的高次扩张。例如把一个四次扩张直接添进去（没有中间二次域）就未必可作。

再举一个经典结论（不必记证明）：正 $n$ 边形可作当且仅当

\[
n=2^a\,p_1p_2\cdots p_r,
\]

其中 $a\ge0$，$p_i$ 是互不相同的 \textbf{Fermat 素数}（形如 $2^{2^k}+1$ 的素数，如 $3,5,17,257,65537$）。例如正五边形可作、正七边形不可作。

\subsection{本节的勘误与提醒}


\begin{itemize}
\item \textbf{塔律的写法与方向。} 是 $[H:F]=[H:K][K:F]$，不要把两个因子写反；推论是“[K:F] 整除 [H:F]”，不是反过来。
\item \textbf{“可作数”不等于“代数数”。} 超越数一定不可作（如 $\pi$），但代数数也未必可作。
\item \textbf{倍立方是“体积变成 $2$ 倍”，边长是 $\sqrt[3]{2}$}，不要与“边长为 $2$”混淆。
\item \textbf{化圆为方的关键是 $\pi$ 超越，不是 $\pi$ 不是有理数。} 需要说明 $\sqrt\pi$ 代数与 $\pi$ 代数等价。
\item \textbf{每一步必须有限次操作。} 尺规作图不允许“取极限”，所以域塔只有有限层。
\item 课堂转写把“倍立方”“三等分角”“化圆为方”识别成了“被立方体”“三种分解”“化缘为方”之类的词，以本节术语为准。
\end{itemize}

\section{分裂域：存在性与唯一性}


第 1 节用 $F[x]/(f)$ 添了一个根，但一个多项式可能有多个、甚至复的根本不在这个域里。这一节定义\textbf{分裂域}——让整个多项式一次分裂完的“最小”扩域——并证明它\textbf{存在}、在\textbf{同构意义下唯一}。有限域的存在性正要用到这条定理：$\mathbb F_{p^d}$ 就是 $x^{p^d}-x$ 的分裂域。

\textbf{本节用到的记号：} $F,E,K$ 为域；$f\in F[x]$ 为次数 $\ge1$ 的多项式（总可除以首项系数，不妨设 $f$ 首一，常数倍不影响根和分裂）；$\alpha_i$ 为 $f$ 在扩域中的根；$F(\alpha_1,\dots,\alpha_n)$ 为 $F$ 与这些根生成的最小子域；$F[x]/(f)$ 为商环；$\cong$ 表示域同构。

\subsection{分裂与分裂域}


设 $f\in F[x]$ 是次数 $n\ge1$ 的多项式。若扩域 $E\supseteq F$ 中有 $\alpha_1,\dots,\alpha_n\in E$ 及 $c\in F^\times$ 使

\[
f(x)=c\,(x-\alpha_1)(x-\alpha_2)\cdots(x-\alpha_n),
\]

就说 $f$ 在 $E$ 上\textbf{完全分裂}，并把 $\alpha_1,\dots,\alpha_n$ 称为 $f$ 在 $E$ 中的根（重根按重数记）。称 $E$ 是 $f$ 的\textbf{分裂域}，若

\begin{itemize}
\item $f$ 在 $E$ 上完全分裂；
\item $E$ 由 $F$ 和 $f$ 的全部根生成，即 $E=F(\alpha_1,\dots,\alpha_n)$。
\end{itemize}


这里的第 2 条正是“最小”的准确写法：\textbf{分裂域只添了必要的根，没有多余元素。} 只要求“$f$ 在 $E$ 上分裂”是不够的——把 $E$ 扩大，$f$ 仍然分裂，所以必须要求“由根生成”。

\textbf{根与一次因子。} 设 $E\supseteq F$、$\alpha\in E$。把 $f$ 看成 $E[x]$ 中的多项式，作带余除法：

\[
f(x)=(x-\alpha)q(x)+r,\qquad r\in E\ \text{为常数}.
\]

代入 $x=\alpha$ 得 $r=f(\alpha)$。因此

\[
f(\alpha)=0\iff (x-\alpha)\mid f(x)\ \text{在}\ E[x]\ \text{中}.
\]

这条小事实反复使用：\textbf{“添一个根”与“多一个一次因子”是同一件事。}

\subsection{存在性：逐步添根}


\textbf{定理。} 对任意域 $F$ 与任意次数 $n\ge1$ 的多项式 $f\in F[x]$，$f$ 的分裂域存在；并且任一分裂域 $E$ 满足

\[
[E:F]\le n!.
\]

\textbf{证明（对 $n=\deg f$ 归纳）。}

\textbf{基础情形 $n=1$。} 设 $f=c(x-\alpha)$，$\alpha\in F$。此时 $f$ 已在 $F$ 中分裂，$E=F$，$[E:F]=1=1!$。

\textbf{归纳步。} 设 $n\ge2$，并假设次数小于 $n$ 的多项式都有分裂域且次数上界成立。把 $f$ 在 $F[x]$ 中分解成不可约多项式之积

\[
f=c\,g_1g_2\cdots g_r .
\]

若每个 $g_i$ 都是一次，则 $f$ 已在 $F$ 上分裂，取 $E=F$ 即可（此时 $[E:F]=1\le n!$）。否则存在某个 $q=g_i$ 满足 $d=\deg q\ge2$。令

\[
F_1=F[x]/(q),
\]

则 $F_1$ 是域、$[F_1:F]=d\le n$，且 $\alpha=\bar x\in F_1$ 满足 $q(\alpha)=0$。因为 $q\mid f$，所以 $f(\alpha)=0$，于是由“根 $\iff$ 一次因子”，在 $F_1[x]$ 中

\[
f(x)=(x-\alpha)\,f_1(x),\qquad \deg f_1=n-1.
\]

把 $f_1$ 看成 $F_1$ 上的次数 $n-1$ 多项式，由归纳假设，$f_1$ 在 $F_1$ 上存在分裂域 $E$，且

\[
[E:F_1]\le (n-1)!.
\]

因为 $E\supseteq F_1\supseteq F$，且 $E$ 中已经有 $f_1$ 的所有根、也已经有 $\alpha$（$F_1$ 的元素），所以 $f=(x-\alpha)f_1$ 在 $E$ 上完全分裂；又 $E$ 由 $F_1$ 与 $f_1$ 的根生成，而 $F_1=F(\alpha)$，所以

\[
E=F(\alpha,\ \text{$f_1$ 的全部根})=F(\text{$f$ 的全部根}),
\]

即 $E$ 是 $f$ 的分裂域。由塔律，

\[
[E:F]=[E:F_1][F_1:F]\le (n-1)!\cdot n=n!.
\]

归纳完成。$\square$

\textbf{证明在做什么。} 每一步只添一个根，把多项式的次数降 $1$；重复至多 $n$ 次，就把所有根添完了。上界也顺势得到：每步的次数至多等于当时多项式的次数，乘积不超过 $n!$。

\subsection{例子}


\textbf{例一：$f=x^2-2$ 在 $\mathbb Q$ 上。} 两个根 $\pm\sqrt2$。分裂域

\[
E=\mathbb Q(\sqrt2),\qquad [E:\mathbb Q]=2=2! .
\]

负根不用单独添：$-\sqrt2=(-1)\cdot\sqrt2$，而 $-1\in\mathbb Q$。

\textbf{例二：$f=x^3-2$ 在 $\mathbb Q$ 上。} 三个根是

\[
\sqrt[3]{2},\qquad \sqrt[3]{2}\,\omega,\qquad \sqrt[3]{2}\,\omega^2,
\qquad \omega=e^{2\pi i/3}=\frac{-1+\sqrt{-3}}{2}.
\]

只添实根 $\sqrt[3]{2}$ 不够：另外两个是复根，而 $\mathbb Q(\sqrt[3]{2})\subseteq\mathbb R$。于是

\[
E=\mathbb Q(\sqrt[3]{2},\omega).
\]

计算次数：$[\mathbb Q(\sqrt[3]{2}):\mathbb Q]=3$（$x^3-2$ 不可约）；在 $\mathbb Q(\sqrt[3]{2})$ 上，$\omega$ 满足 $x^2+x+1=0$，而这个二次多项式在实域 $\mathbb Q(\sqrt[3]{2})$ 中没有根，所以不可约，$[\mathbb Q(\sqrt[3]{2},\omega):\mathbb Q(\sqrt[3]{2})]=2$。由塔律

\[
[E:\mathbb Q]=3\cdot2=6=3! .
\]

这个例子达到了上界 $n!$，说明上界一般不能改进。

\textbf{例三：$f=x^2-1$ 在 $\mathbb Q$ 上。} 根 $\pm1$ 已经都在 $\mathbb Q$ 中，所以 $E=\mathbb Q$，$[E:\mathbb Q]=1<2!$。上界不是等号。

\textbf{例四（前向引用）：$f=x^{p^d}-x$ 在 $\mathbb F_p$ 上。} 它的分裂域就是 $\mathbb F_{p^d}$，这一点下一节用来证明有限域存在。这里先记住：它是“每个元素都满足 $x^{p^d}=x$”的域。

\subsection{唯一性：同构意义下唯一}


存在性之外，还要说明分裂域不依赖于“先添哪个根”的选择。先准备一个延拓引理。

\textbf{引理（同构延拓）。} 设 $\sigma:F\to F'$ 是域同构，$f\in F[x]$ 不可约，记 $f^\sigma\in F'[x]$ 为把 $f$ 的系数逐个用 $\sigma$ 作用得到的多项式。设 $E/F$、$E'/F'$ 是扩域，$\alpha\in E$ 是 $f$ 的根，$\alpha'\in E'$ 是 $f^\sigma$ 的根。则存在域同构

\[
\tilde\sigma:F(\alpha)\longrightarrow F'(\alpha')
\]

使得 $\tilde\sigma|_F=\sigma$ 且 $\tilde\sigma(\alpha)=\alpha'$。

\textbf{证明。} 定义求值映射

\[
\psi:F[x]\longrightarrow F'(\alpha'),\qquad g(x)\longmapsto \sigma(g)(\alpha'),
\]

其中 $\sigma(g)$ 表示把 $g$ 的系数用 $\sigma$ 作用。$\psi$ 是环同态（$\sigma$ 是同构），且 $\psi(1)=1\ne0$。因为 $f$ 不可约，$(f)$ 是 $F[x]$ 的极大理想；又 $f\in\ker\psi$（因为 $\psi(f)=f^\sigma(\alpha')=0$）且 $\psi(1)=1\ne0$，所以 $\ker\psi$ 是包含 $(f)$ 的真理想，只能是 $(f)$ 本身。于是

\[
F[x]/(f)\ \cong\ \operatorname{im}\psi .
\]

左边是域（$f$ 不可约），所以 $\operatorname{im}\psi$ 是域，并且包含 $F'$（常数多项式的像）与 $\alpha'$（$x$ 的像）。由 $F'(\alpha')$ 的最小性，$\operatorname{im}\psi=F'(\alpha')$。

另一方面，上一节已证 $F(\alpha)\cong F[x]/(f)$，同构把 $\bar x$ 送到 $\alpha$。复合

\[
F(\alpha)\xrightarrow{\ \cong\ }F[x]/(f)\xrightarrow{\ \cong\ }F'(\alpha')
\]

就得到 $\tilde\sigma$；检查两个同构在 $F$ 上都等于 $\sigma$、在 $\alpha$ 上都对应 $\alpha'$，故 $\tilde\sigma$ 满足要求。$\square$

\textbf{定理（分裂域唯一）。} 设 $E,E'$ 都是 $f\in F[x]$ 在 $F$ 上的分裂域。则存在域同构 $\sigma:E\to E'$ 使 $\sigma|_F=\operatorname{id}_F$。

\textbf{证明（对 $n=\deg f$ 归纳）。} $n=1$ 时 $E=E'=F$，取恒等映射。设 $n\ge2$。

若 $f$ 已在 $F$ 上分裂，则 $E=E'=F$，结论显然。否则取 $f$ 的一个次数 $\ge2$ 的不可约因子 $q\in F[x]$。因为 $q\mid f$ 且 $f$ 在 $E$、$E'$ 中都分裂，$q$ 在 $E$、$E'$ 中都有根，记为 $\alpha$、$\alpha'$。由延拓引理（取 $\sigma=\operatorname{id}_F$），存在同构

\[
\sigma_1:F(\alpha)\longrightarrow F(\alpha'),\qquad \sigma_1(\alpha)=\alpha',
\]

且 $\sigma_1$ 在 $F$ 上是恒等。

在 $F(\alpha)[x]$ 中写

\[
f(x)=(x-\alpha)\,f_1(x),\qquad \deg f_1=n-1.
\]

因为 $E$ 由 $F(\alpha)$ 与 $f_1$ 的其余根生成，$E$ 是 $f_1$ 在 $F(\alpha)$ 上的分裂域；同理（把 $f$ 的系数与 $\sigma_1$ 作用）$E'$ 是 $\sigma_1(f_1)$ 在 $F(\alpha')$ 上的分裂域。由归纳假设，存在同构 $\sigma:E\to E'$ 使 $\sigma|_{F(\alpha)}=\sigma_1$。特别地 $\sigma|_F=\operatorname{id}_F$。$\square$

\textbf{唯一性要读准的三件事。}

\begin{itemize}
\item \textbf{“唯一”是同构意义下的唯一，不是集合相等。} 两个分裂域可能由不同的元素、甚至不同的构造得到；它们之间保持 $F$ 不动的同构存在，就已经是能做到的最好结果。
\item \textbf{如果 $E,E'$ 都放在同一个更大的域 $K$ 里，则它们作为子域相等。} 因为两者都恰好是 $K$ 中由 $F$ 与全部根生成的子域。课程讲的“唯一性”通常指这种情形。
\item \textbf{分裂域没有多余元素。} 若 $E$ 分裂 $f$ 且 $E$ 真包含于“由根生成的子域”，那 $E$ 不是分裂域；若 $E$ 严格大于由根生成的子域，多出来的元素应删去。
\end{itemize}


\subsection{补充：唯一性定理的常见用法}


\textbf{用法一：计算分裂域的次数。} 若已知 $E=F(\alpha_1,\dots,\alpha_n)$，可逐层用塔律；唯一性保证这个次数不依赖取根顺序。

\textbf{用法二：把“形式构造”与“具体扩域”对照。} 存在性证明给出一个由商环逐步搭起来的域；具体问题中往往先看到某个大域里的根，唯一性保证两者同构。

\textbf{用法三：有限域唯一性的证明模板。} 下一节将证明：任何 $p^d$ 元域都恰好是 $x^{p^d}-x$ 在 $\mathbb F_p$ 上的分裂域，故由本节的唯一性，$p^d$ 元域在同构意义下唯一。

\subsection{本节的勘误与提醒}


\begin{itemize}
\item \textbf{分裂域必须“最小”。} 只说“$f$ 在 $E$ 上分裂”不构成定义；例如 $\mathbb C$ 分隔 $x^2-2$，但 $\mathbb C$ 不是它的分裂域。
\item \textbf{分裂域的存在性证明用的是“添根 + 归纳”。} 每一步添的是形式根 $\bar x$，不是先假定根存在于某个已知大域中；这是存在性证明的关键。
\item \textbf{上界 $n!$ 不是等式。} $x^2-1$ 给出 $1<2$；$x^3-2$ 才达到 $6$。
\item \textbf{“同构意义下唯一”不能省。} 课堂转写里“存在唯一分裂域”应补成“存在分裂域，且在 $F$-同构意义下唯一”。
\item \textbf{唯一性证明没有用到特征 $0$}，对任意域（包括有限域）都成立；这是下一节能用的原因。
\item 转写中“分裂域的唯一性”与“存在性”经常混在一起说，读时按本节两个定理区分。
\end{itemize}

\section{有限域：存在唯一性与不可约多项式的计数}


Lec10 已经证明：有限域的阶必为素数幂 $p^d$。这一节补上反方向——\textbf{每个素数幂都存在有限域，且在同构意义下唯一}——并回答一个看似无关的问题：\textbf{每个次数 $d$ 是否一定存在 $\mathbb F_p$ 上的 $d$ 次不可约多项式？} 答案（存在）正是手写稿第 3 页上半页的内容；它是下一节“构造有限域并做因式分解”的前提。

\textbf{本节用到的记号：} $p$ 为素数，$q=p^r$ 为素数幂（字母 $q$ 在因式分解一节还会出现，含义相同）；$\mathbb F_p=\mathbb Z/p\mathbb Z$；$\mathbb F_q$ 表示 $q$ 元有限域；$f'$ 为形式导数；$\Phi_d(x)$ 表示“$\mathbb F_q$ 上全体 $d$ 次首一不可约多项式之积”（手写稿写成 $\varphi_d$，$q=p$ 时就是 $\mathbb F_p$ 上的情形）；$N_d=N_d(q)$ 表示 $\mathbb F_q$ 上 $d$ 次首一不可约多项式的\textbf{个数}（$q=p$ 时即手写稿里的数）；$\mu$ 为 Möbius 函数。

\subsection{有限域的大小（回顾）}


设 $E$ 是有限域。$\operatorname{char}(E)$ 不能是 $0$（否则 $E$ 含 $\mathbb Q$，已经无限），故特征为素数 $p$，且 $E$ 含素子域 $\mathbb F_p$。于是 $E$ 是 $\mathbb F_p$ 上的有限维线性空间，设 $[E:\mathbb F_p]=d$，则

\[
|E|=p^d .
\]

\textbf{有限域的大小必为素数幂。} 反过来的存在唯一性就是本节主定理。

\subsection{Frobenius 引理：特征 $p$ 的“加法变简单”}


\textbf{引理（Frobenius 引理）。} 设 $K$ 是特征 $p$ 的域。映射

\[
\sigma_p:K\longrightarrow K,\qquad a\longmapsto a^p
\]

是域同态，且是单射；它的不动域 $\{a\in K:a^p=a\}$ 恰为素子域 $\mathbb F_p$。若 $K$ 还是\textbf{有限域}，则 $\sigma_p$ 是双射，从而是域自同构。更一般地，$a\mapsto a^{q}$（$q=p^r$）也是域同态，称为 \textbf{$q$-Frobenius 映射}（$K$ 有限时同样是自同构）。

\textbf{证明。} $(ab)^p=a^pb^p$、$1^p=1$ 是显然的。加法性来自二项式定理：

\[
(a+b)^p=\sum_{k=0}^{p}\binom{p}{k}a^k b^{p-k}=a^p+b^p,
\]

因为中间的二项式系数 $\binom{p}{k}$（$1\le k\le p-1$）都是 $p$ 的倍数，而在特征 $p$ 的域中 $p=0$，这些项都消失。又 $\sigma_p$ 是域同态，由“域同态非零则单射”（第 1 节），它是单射。当 $K$ 是有限域时，从有限集到自身的单射必是双射，所以此时 $\sigma_p$ 是自同构。

不动域是一个子域：若 $a^p=a$、$b^p=b$，则 $(a\pm b)^p=a^p\pm b^p=a\pm b$、$(ab)^p=ab$、$(a^{-1})^p=(a^p)^{-1}=a^{-1}$。$x^p-x$ 在域中至多 $p$ 个根，而 $\mathbb F_p$ 的每个元素都是它的根，所以不动域恰为 $\mathbb F_p$。多次复合 $\sigma_p^r$ 就得到 $\sigma_q$。$\square$

\textbf{推论（Fermat 小定理的有限域版）。} 若 $E$ 是 $q$ 元有限域，则任意 $a\in E$ 满足 $a^q=a$；任意 $a\ne0$ 满足 $a^{q-1}=1$。

\textbf{证明。} $E^\times$ 是 $q-1$ 阶乘法群，由 Lagrange 定理 $a^{q-1}=1$；再乘 $a$ 得 $a^q=a$；$a=0$ 显然。$\square$

\subsection{存在性：$x^{p^d}-x$ 的分裂域}


\textbf{定理。} 对任意素数 $p$ 与正整数 $d$，存在 $p^d$ 元有限域。

\textbf{证明。} 考虑多项式

\[
f(x)=x^{p^d}-x\ \in\ \mathbb F_p[x].
\]

由上一节，$f$ 在 $\mathbb F_p$ 上有分裂域 $K$，且 $[K:\mathbb F_p]\le(p^d)!$，所以 $K$ 是有限域。下面证明 $K$ 恰有 $p^d$ 个元素。

\textbf{第一步：$f$ 没有重根。} 形式导数为

\[
f'(x)=p^d x^{p^d-1}-1=-1\ne0,
\]

因为 $p^d=0$ 于特征 $p$。一般地，多项式有重根当且仅当它与导数有公因式（若 $\alpha$ 是重根，则 $x-\alpha$ 同时整除 $f$ 与 $f'$）；这里 $f'=-1$ 与 $f$ 互质，故 $f$ 无重根。于是 $f$ 在 $K$ 中恰有 $p^d$ 个不同的根。

\textbf{第二步：根集是子域。} 记

\[
R=\{a\in K: a^{p^d}=a\}.
\]

它恰是“$f$ 的全部根”的集合。由 Frobenius 引理，$\sigma_{p^d}:a\mapsto a^{p^d}$ 是域同态，于是 $R$ 是 $\sigma_{p^d}$ 的不动点集，因而是 $K$ 的子域：$0,1\in R$，且对 $a,b\in R$ 有 $a\pm b,ab\in R$，$a\ne0$ 时 $a^{-1}\in R$。

也可以直接验证加法封闭性（课堂写法）：由 Frobenius，$(a+b)^{p^d}=a^{p^d}+b^{p^d}=a+b$；乘法同理。

\textbf{第三步：$K=R$。} 一方面 $R\subseteq K$。另一方面，$K$ 是 $f$ 的分裂域，即由 $\mathbb F_p$ 与 $f$ 的全部根生成；而全部根都在 $R$ 中，$R$ 又是包含 $\mathbb F_p$ 的域，所以由生成性 $K\subseteq R$。因此 $K=R$，$|K|=|R|=p^d$。$\square$

\textbf{这个证明的要点。} 与其“凭空造”一个 $p^d$ 元域，不如拿一个\textbf{其根恰好应有 $p^d$ 个且能构成域}的多项式：$x^{p^d}-x$。根集对四则运算封闭（Frobenius 的功劳），而它又自动是分裂域（因为分裂域恰好由这些根生成）。

\textbf{记号。} 把 $p^d$ 元有限域记为 $\mathbb F_{p^d}$（也写作 $\operatorname{GF}(p^d)$）。由本定理它是分裂域，故存在。

\subsection{唯一性：任意两个 $p^d$ 元域同构}


\textbf{定理。} 任意两个 $p^d$ 元有限域都同构。

\textbf{证明。} 设 $E$ 是 $p^d$ 元域。它特征为 $p$，含素子域 $\mathbb F_p$，且 $[E:\mathbb F_p]=d$。由 Fermat 小定理的有限域版，每个 $a\in E$ 都满足 $a^{p^d}=a$，即 $E$ 中每个元素都是 $f=x^{p^d}-x$ 的根。于是 $f$ 在 $E$ 上分裂，$E$ 由 $\mathbb F_p$ 与 $f$ 的全部根生成（因为所有元素都是根），所以 $E$ 是 $f$ 在 $\mathbb F_p$ 上的分裂域。由分裂域在 $F$-同构意义下唯一（上一节），$E$ 与上一小节构造的 $\mathbb F_{p^d}$ 同构。$\square$

\textbf{注意唯一性的范围。} 唯一的是\textbf{同构类}：两个 $p^d$ 元域之间总有保持 $\mathbb F_p$ 不动的同构，但它们作为集合未必相等（例如同一抽象域的不同具体模型）。

\subsection{推广：$\mathbb F_q$ 上的 $x^{q^n}-x$}


把上面的论证逐字替换 $p\mapsto q$、$d\mapsto n$，得到：

\textbf{定理（一般 $q$ 版）。} 设 $q=p^r$。对任意 $n\ge1$，$x^{q^n}-x$ 在 $\mathbb F_q$ 上的分裂域是 $q^n$ 元域 $\mathbb F_{q^n}$；它是 $\mathbb F_q$ 的唯一（同构意义下）$n$ 次扩域。

证明完全一样：$f'=-1$ 无重根；根集是 Frobenius $\sigma_{q^n}$ 的不动域；分裂域由根生成。以后因式分解一节就在一般的 $\mathbb F_q$ 上工作，取 $q=p$ 就回到手写稿的写法。

\textbf{顺带：子域结构。} 对固定的 $p$，$\mathbb F_{p^m}\subseteq\mathbb F_{p^n}$ 当且仅当 $m\mid n$；$\mathbb F_{p^n}$ 的每个子域都是某个 $\mathbb F_{p^m}$（$m\mid n$）。证明：若 $m\mid n$，则 $x^{p^m}=x\Rightarrow x^{p^n}=x$，所以 $\mathbb F_{p^m}$（$x^{p^m}-x$ 的根集）包含于 $\mathbb F_{p^n}$；反之若 $\mathbb F_{p^m}\subseteq\mathbb F_{p^n}$，塔律给出 $n=[\mathbb F_{p^n}:\mathbb F_p]=[\mathbb F_{p^n}:\mathbb F_{p^m}]\cdot m$，故 $m\mid n$。

\subsection{核心分解：$x^{p^d}-x$ 与不可约多项式}


这一小节回答“$d$ 次不可约多项式是否存在、有多少个”。

\textbf{定理（分解恒等式）。} 设 $q=p^r$，则

\[
x^{q^d}-x=\prod_{e\mid d}\Phi_e(x),
\]

其中 $\Phi_e(x)$ 是 $\mathbb F_q$ 上全体 $e$ 次首一不可约多项式之积。（本讲开头手写稿主要写 $q=p$ 的情形，即 $x^{p^d}-x=\prod_{e\mid d}\Phi_e$。）

\textbf{证明。} 两边都是首一多项式。我们证明“每个首一不可约多项式在两边出现的次数相同”。

先说明 $x^{q^d}-x$ 无重根：它的导数是常数 $-1$，与 $x^{q^d}-x$ 本身互质。所以左边每个不可约因子只出现一次。

于是只需证明：**$\mathbb F_q$ 上的首一不可约多项式 $g$ 满足

\[
g\mid x^{q^d}-x\iff \deg g\mid d.
\]

\textbf{（$\Rightarrow$）} 设 $g\mid x^{q^d}-x$，取 $g$ 的一个根 $\alpha$（在某个扩域中），则 $\alpha^{q^d}=\alpha$，即 $\alpha$ 属于 $x^{q^d}-x$ 的根集，也就是 $\mathbb F_{q^d}$。于是 $\mathbb F_q(\alpha)\subseteq\mathbb F_{q^d}$。因为 $g$ 不可约，$\mathbb F_q(\alpha)\cong\mathbb F_q[x]/(g)$，故 $[\mathbb F_q(\alpha):\mathbb F_q]=\deg g$，由塔律 $\deg g\mid d$。

\textbf{（$\Leftarrow$）} 设 $e=\deg g\mid d$。取 $g$ 的一个根 $\alpha$。因为 $g$ 是 $\mathbb F_q$ 上以 $\alpha$ 为根的首一不可约多项式，所以 $g$ 就是 $\alpha$ 的极小多项式；而极小多项式整除任何以 $\alpha$ 为根的 $\mathbb F_q$ 系数多项式，故 $g\mid x^{q^d}-x$。

（也可以把理由拆开看：由 $\mathbb F_q(\alpha)\cong\mathbb F_{q^e}\subseteq\mathbb F_{q^d}$（子域判据）得 $\alpha^{q^d}=\alpha$，即 $\alpha$ 是 $x^{q^d}-x$ 的根；再用极小多项式性质即可。）

综上，左边的不可约因子恰好是“次数整除 $d$ 的首一不可约多项式”，每个一次，正是右边。$\square$

\textbf{这个恒等式在说什么。} 它把 $x^{q^d}-x$ 按“不可约因子的次数”\textbf{分层}：所有 $e$ 次不可约多项式（$e\mid d$）之积，恰好等于 $x^{q^d}-x$。这就是手写稿第 3 页上半页的核心。

\textbf{补充引理（gcd 形式，课堂提到）。} 对 $a,b\ge1$，

\[
\gcd\bigl(x^{q^a}-x,\ x^{q^b}-x\bigr)=x^{q^{\gcd(a,b)}}-x .
\]

\textbf{证明思路。} 不妨设 $a\ge b$。在模 $x^{q^b}-x$ 下 $x^{q^b}\equiv x$，两边取 $q^{a-b}$ 次幂得 $x^{q^a}\equiv x^{q^{a-b}}$，所以 $\gcd(x^{q^a}-x,x^{q^b}-x)=\gcd(x^{q^{a-b}}-x,x^{q^b}-x)$；对指数做辗转相除即得。$\square$

\subsection{计数：$N_d$ 与 Möbius 公式}


设 $N_e=N_e(q)$ 为 $\mathbb F_q$ 上 $e$ 次首一不可约多项式的个数。由 $\Phi_e$ 的定义，$\deg\Phi_e=e\,N_e$。对分解恒等式两边比较次数：

\[
q^d=\sum_{e\mid d} e\,N_e .
\qquad(\ast)
\]

这是关于 $N_d$ 的\textbf{递归公式}：

\[
d\,N_d=q^d-\sum_{e\mid d,\ e<d} e\,N_e .
\]

\textbf{小数值。} 取 $q=p$：

\begin{itemize}
\item $N_1=p$（一次首一多项式 $x-a$，$a\in\mathbb F_p$）；
\item $N_2=(p^2-p)/2$；
\item $N_3=(p^3-p)/3$；
\item $N_4=(p^4-p^2)/4$；
\item $N_6=(p^6-p^3-p^2+p)/6$。
\end{itemize}


这些正是手写稿上 $\deg\Phi_1=p$、$\deg\Phi_2=p^2-p$、$\deg\Phi_3=p^3-p$、$\deg\Phi_6=p^6-p^3-p^2+p$ 的除法版本：$\deg\Phi_d=d\,N_d$。

具体数表（首一不可约多项式个数）：

\begin{itemize}
\item $p=2$：$N_1,\dots,N_8=2,1,2,3,6,9,18,30$；
\item $p=3$：$3,3,8,18,48,116,312,810$。
\end{itemize}


\textbf{Möbius 反演（进阶，手写稿没写，可作补充）。} 定义 Möbius 函数

\[
\mu(n)=
\begin{cases}
1, & n=1,\\
0, & n\ \text{含平方因子},\\
(-1)^k, & n\ \text{是}\ k\ \text{个不同素数之积}.
\end{cases}
\]

由 $(\ast)$ 对 $d$ 做 Möbius 反演得

\[
N_d=\frac1d\sum_{e\mid d}\mu\!\left(\frac de\right)q^e .
\]

验证：$q=2,d=4$ 时 $N_4=\frac14(\mu(1)2^4+\mu(2)2^2+\mu(4)2^1)=\frac14(16-4+0)=3$。

\subsection{每个次数都有不可约多项式}


有了 $(\ast)$ 与上面的“$\deg\Phi_e\le q^e$”（因为 $\Phi_e\mid x^{q^e}-x$），可以证明：

\textbf{定理。} 对任意素数幂 $q$ 与任意 $d\ge1$，存在 $d$ 次不可约多项式。

\textbf{证明。} 由分解恒等式，

\[
\deg\Phi_d=q^d-\sum_{e\mid d,\ e<d}\deg\Phi_e
\ \ge\ q^d-\sum_{e\mid d,\ e<d}q^e .
\]

真正的真因子 $e$ 满足 $e\le d/2$（$d\ge2$ 时）。于是

\[
\sum_{e\mid d,\ e<d}q^e
\le q+q^2+\cdots+q^{\lfloor d/2\rfloor}
=\frac{q^{\lfloor d/2\rfloor+1}-q}{q-1}
<q^{d/2+1}
\le q^d,
\]

最后一步用了 $d\ge2$。因此 $\deg\Phi_d>0$，即至少有一个 $d$ 次不可约多项式。$d=1$ 时 $\deg\Phi_1=q>0$ 显然。$\square$

\textbf{这条定理的用途。} 它说明构造 $\mathbb F_{q^d}$ 有两条路：

\begin{itemize}
\item \textbf{分裂域路线}：取 $x^{q^d}-x$ 的分裂域（存在性证明用的）；
\item \textbf{商环路线}：任取一个 $d$ 次不可约 $g\in\mathbb F_q[x]$，则 $\mathbb F_q[x]/(g)$ 是 $q^d$ 元域，且 $\mathbb F_{q^d}\cong\mathbb F_q[x]/(g)$。
\end{itemize}


两条路殊途同归（唯一性保证）。实际计算用第二条：只要找到 $g$ 就能写出域的元素与运算。

\subsection{例子}


\textbf{例一：$p=2$ 的逐层分解。}

\[
\begin{aligned}
x^2-x&=\Phi_1=x(x+1),\\
x^4-x&=\Phi_1\Phi_2=x(x+1)(x^2+x+1),\\
x^8-x&=\Phi_1\Phi_3=x(x+1)(x^3+x+1)(x^3+x^2+1),\\
x^{16}-x&=\Phi_1\Phi_2\Phi_4=x(x+1)(x^2+x+1)(x^4+x+1)(x^4+x^3+1)(x^4+x^3+x^2+x+1).
\end{aligned}
\]

这里 $d=4$ 的因子只有 $1,2,4$，没有 $3$，所以 $x^{16}-x$ 中不出现三次因子；$N_4=3$ 对应最后三个四次不可约多项式。

\textbf{例二：$p=3$ 的逐层分解。}

\[
x^3-x=x(x-1)(x+1),\qquad
x^9-x=x(x-1)(x+1)(x^2+1)(x^2+x-1)(x^2-x-1).
\]

第二个式子中 $N_2=3$，三个二次不可约多项式是 $x^2+1$、$x^2+x-1$、$x^2-x-1$。

\textbf{例三：$\mathbb F_{16}$ 的元素按极小多项式分组。} 取 $\mathbb F_{16}=\mathbb F_2[x]/(x^4+x+1)$，16 个元素按“在 $\mathbb F_2$ 上的极小多项式”分类：

\begin{itemize}
\item $0$ 对应 $x$，$1$ 对应 $x+1$：2 个元素；
\item $x^2+x+1$ 的 2 个根：2 个元素；
\item 三个四次不可约多项式各 4 个根：12 个元素。
\end{itemize}


合计 $2+2+12=16$，与 $\sum_{e\mid4}eN_e=1\cdot2+2\cdot1+4\cdot3=16$ 一致。

\textbf{例四：具体构造。}

\[
\mathbb F_4=\mathbb F_2[x]/(x^2+x+1),\qquad
\mathbb F_8=\mathbb F_2[x]/(x^3+x+1),\qquad
\mathbb F_9=\mathbb F_3[x]/(x^2+1).
\]

在 $\mathbb F_9$ 中把 $x$ 的像记为 $i$（$i^2=-1=2$），则元素为 $a+bi$（$a,b\in\mathbb F_3$），按“$i^2=2$ 降次”做乘法。

\subsection{本节的勘误与提醒（对应手写稿第 3 页上半页）}


\begin{itemize}
\item \textbf{手写稿的 $\varphi_d$ 是多项式，不是欧拉函数。} 定义为“全体 $d$ 次首一不可约多项式之积”，所以 $\deg\varphi_d=d\cdot N_d$；它和数论里的欧拉函数 $\varphi(n)$ 只是记号撞车，含义无关。
\item \textbf{手写稿“$x^{p^6}-x=\varphi_6+\varphi_3+\varphi_2+\varphi_1$”中的加号是笔误，应为乘积。} 依据是分解恒等式；比较次数也对不上加号。
\item \textbf{不等式的下标是 $d'\mid d$ 且 $d'\ne d$。} 手写稿的 $\sum$ 下方同时写了 $d'\mid d$ 与 $d'\ne d$（即“$d$ 的真因子”，与 $d'<d$ 等价）；读的时候注意别漏掉 $d'\mid d$，否则会变成无穷和。
\item \textbf{“每个次数都有不可约多项式”的证明不是构造性的。} 它只证明存在；要真正找到一个，可以枚举并检验（或用下一节的算法思想）。
\item \textbf{$f\mid x^{q^d}-x\iff\deg f\mid d$ 的方向不能反。} 例如 $x+1$（次数 $1$）整除任何 $x^{q^d}-x$；而四次不可约多项式只整除满足 $4\mid d$ 的那些。
\item \textbf{$x^{q^d}-x$ 无重根。} 这是“分解恒等式”与下一节“按次数分组”算法都依赖的事实；来源是导数 $-1$。
\item \textbf{有限域大小必须是素数幂}：不存在 $6$ 元域。因为若 $|E|=n$，则 $n=p^{[E:\mathbb F_p]}$。
\item \textbf{不同不可约多项式给出的域同构，但元素对应关系不同。} 例如 $\mathbb F_4=\mathbb F_2[x]/(x^2+x+1)$ 中 $x$ 的像是“$\omega$”，换一个二次不可约多项式（$\mathbb F_2$ 上只有这一个）没有区别；$\mathbb F_9$ 上不同的二次不可约多项式会给出不同的生成元对应。
\item 课堂转写把 Frobenius 的“$(a+b)^{p^d}=a^{p^d}+b^{p^d}$”说成“二项式展开中间项消失了”，这是对的；但注意\textbf{只在特征 $p$ 下成立}，一般环里不成立。
\end{itemize}

\section{有限域上的多项式因式分解}


这是本讲的最后一个主题，也是手写稿第 3 页下半页的内容：\textbf{给定有限域 $\mathbb F_q$ 上的多项式 $f$，怎样高效地把 $f$ 分解成不可约多项式之积。} 在整数里分解质因数被认为是很难的问题（没有量子计算机时）；但在有限域的多项式环里，分解可以做到多项式时间——这正是这页三个算法要说明的事。

\textbf{本节用到的记号：} $q=p^r$ 为素数幂，$\mathbb F_q$ 为有限域；$f\in\mathbb F_q[x]$ 首一、次数 $\ge1$；$f'$ 为\textbf{形式导数}；$\gcd$ 表示首一最大公因式（在 $\mathbb F_q[x]$ 中算）；称 $f$ \textbf{无平方因子}，若它没有重复的不可约因子（即每个不可约因子只出现一次）。

\subsection{问题与递归策略}


\textbf{输入：} 首一多项式 $f\in\mathbb F_q[x]$。
\textbf{输出：} 不可约多项式的乘积分解

\[
f=f_1^{e_1}f_2^{e_2}\cdots f_s^{e_s},
\]

其中 $f_i$ 首一不可约、互不相同，$e_i\ge1$。

\textbf{递归策略。} 只要能找到一个\textbf{非平凡因子} $g$（$1<\deg g<\deg f$，$g\mid f$），就可以递归分解 $g$ 与 $f/g$。所以下面三个算法都只负责回答一个问题：

\begin{quote}
怎样从 $f$ 里“高效地切出一刀”？
\end{quote}


整条流水线是：

\begin{itemize}
\item \textbf{去平方因子}：把重复的因子与重数处理好，剩下无平方因子的部分；
\item \textbf{按次数分组}（distinct-degree）：把不可约因子按次数 $1,2,3,\dots$ 分成一堆一堆；
\item \textbf{同次数拆分}（equal-degree，Cantor–Zassenhaus）：把一堆次数相同、但互不相同的不可约因子逐个分开。
\end{itemize}


三步都做完，就得到完整分解。

\subsection{形式导数与重根}


设

\[
f(x)=a_nx^n+a_{n-1}x^{n-1}+\cdots+a_1x+a_0 .
\]

定义\textbf{形式导数}

\[
f'(x)=n a_nx^{n-1}+(n-1)a_{n-1}x^{n-2}+\cdots+a_1 .
\]

它满足通常的乘积法则 $(gh)'=g'h+gh'$（直接把两边按定义展开即可验证，不需要极限）。注意系数 $a_i$ 与整数 $n$ 的乘积要理解成在 $\mathbb F_q$ 中相加 $n$ 次；当 $p\mid n$ 时 $na_n=0$。

\textbf{引理（重因子与导数）。} 若 $g^2\mid f$，则 $g\mid f'$。

\textbf{证明。} 写 $f=g^2h$。由乘积法则

\[
f'=2gg'h+g^2h'=g\,(2g'h+gh'),
\]

故 $g\mid f'$。$\square$

\textbf{推论。} 若 $\gcd(f,f')=1$，则 $f$ 无平方因子；反过来，$f$ 无平方因子当且仅当 $\gcd(f,f')=1$。

\textbf{证明。} 写 $f=\prod_i r_i^{e_i}$ 为不可约分解。

先证“有重复因子 $\Rightarrow\gcd\ne1$”。若某个 $e_i\ge2$：在 $r_i$ 处局部写 $f=r_i^{e_i}h$（$r_i\nmid h$），则
\[
f'=e_i r_i^{e_i-1}r_i'h+r_i^{e_i}h'.
\]
\begin{itemize}
\item 若 $p\nmid e_i$，则括号 $e_ir_i'h+r_ih'$ 不被 $r_i$ 整除（这里用到 $r_i\nmid r_i'$：不可约多项式在特征 $p$ 下导数非零，否则它是 $p$ 次幂），故 $r_i^{e_i-1}\mid f'$，与 $r_i^{e_i}\mid f$ 一起给出 $r_i\mid\gcd(f,f')$；
\item 若 $p\mid e_i$，则 $e_ir_i^{e_i-1}r_i'h=0$，$f'=r_i^{e_i}h'$，仍有 $r_i\mid f'$。
\end{itemize}


所以 $\gcd(f,f')\ne1$。

再证“无重复因子 $\Rightarrow\gcd=1$”。设 $f=\prod_i r_i$（互不相同）。若某个不可约 $r_i\mid f'$，则 $f'\equiv0\pmod{r_i}$。但 $f=\prod_j r_j$，由乘积法则
\[
f'\equiv r_i'\prod_{j\ne i}r_j \pmod{r_i},
\]
而 $r_i\nmid r_i'$、$r_i\nmid\prod_{j\ne i}r_j$，矛盾（域上无零因子）。故没有 $r_i$ 整除 $f'$，$\gcd(f,f')=1$。$\square$

\textbf{边界情形 $f'=0$。} 上面推理里用到“$r_i'\ne0$”。什么时候 $f'=0$？逐个单项式看：$(x^{m})'=mx^{m-1}$，所以 $f'=0$ 当且仅当每个非零项的次数都是 $p$ 的倍数，即

\[
f(x)=h(x^p)
\]

对某个 $h\in\mathbb F_q[x]$ 成立。又因为 $\mathbb F_q$ 中 Frobenius 映射 $a\mapsto a^p$ 是双射（第 4 节 Frobenius 引理：有限域上单射即双射），每个系数都有唯一的 $p$ 次方根，于是 $h(x^p)=\bigl(h^{[1/p]}(x)\bigr)^p$，其中 $h^{[1/p]}$ 表示“把每个系数换成它的 $p$ 次方根”。\textbf{所以 $f'=0$ 等价于 $f$ 本身是一个 $p$ 次幂。}

\subsection{第一步：去平方因子（square-free）}


\textbf{目标。} 把 $f$ 写成若干个无平方因子多项式之积（带重数），从而把问题化归到无平方因子的情形。

\textbf{算法。}

\begin{itemize}
\item \textbf{第一步}：计算 $f'$。
\item \textbf{第二步}：若 $f'=0$，则由上面的讨论 $f=h^p$，其中 $h=h^{[1/p]}$。递归分解 $h$，得到 $h=\prod r_i^{e_i}$，则 $f=\prod r_i^{p e_i}$；返回。
\item \textbf{第三步}：否则计算 $c=\gcd(f,f')$，并令 $w=f/c$。
\item 若 $c=1$：$f$ 无平方因子，直接返回（交给第二步）；
\item 若 $c\ne1$：$w=f/c$ \textbf{无平方因子}，但它只含那些重数不被 $p$ 整除的不可约因子各一次；重数被 $p$ 整除的因子全部留在 $c$ 里（详见下一段的说明）。递归对 $c$ 做去平方因子，得到 $c=\prod_m a_m^m$（$a_m$ 无平方因子）；再把 $w$ 的因子按“重数 $+1$”并入，即对每个不可约因子 $r$，它在 $f$ 中的重数等于它在 $c$ 中的重数加上“$r$ 是否出现在 $w$ 中”。
\end{itemize}


\textbf{为什么 $w=f/\gcd(f,f')$ 无平方因子。} 在不可约分解 $f=\prod r_i^{e_i}$ 下，局部看 $r_i$ 的指数：若 $p\nmid e_i$，则 $\gcd$ 中 $r_i$ 的指数是 $e_i-1$，相除后剩 $1$；若 $p\mid e_i$，则 $\gcd$ 中指数是 $e_i$，相除后不出现。两种情形相除后 $r_i$ 的指数都在 $\{0,1\}$，所以 $w$ 无平方因子。

\textbf{例一（重因子）。} 在 $\mathbb F_3$ 上取

\[
f=(x^2+1)^2(x+2)=x^5-x^4-x^3+x^2+x-1 .
\]

计算 $f'=5x^4-4x^3-3x^2+2x+1\equiv -x^4-x^3-x+1$（系数模 $3$）。用辗转相除得

\[
\gcd(f,f')=x^2+1,\qquad f/\gcd(f,f')=x^3-x^2+x-1=(x^2+1)(x+2).
\]

于是 $f=(x^2+1)\cdot\bigl[(x^2+1)(x+2)\bigr]$，重因子 $x^2+1$ 被“降了一次重数”，继续处理剩余部分即可。

\textbf{例二（$f'=0$）。} 在 $\mathbb F_p$ 上取

\[
f=x^{2p}-x^p+3 .
\]

则 $f'=2p\,x^{2p-1}-p\,x^{p-1}=0$。由第 4 节 Frobenius 引理（$(a+b)^p=a^p+b^p$、系数 $c^p=c$ 对 $c\in\mathbb F_p$），

\[
f=(x^2-x+3)^p .
\]

于是 $p$ 次方根为 $x^2-x+3$；递归分解它即可，最后把每个因子的重数乘 $p$。例如 $p=3$ 时 $f=(x^2-x+3)^3=(x^2-x)^3=x^6-x^3$，而 $x^2-x=x(x-1)$，所以 $f=x^3(x-1)^3$。

\textbf{注意。} 去平方因子之后，$f$ 无平方因子，但这\textbf{不表示 $f$ 不可约}。它只是把“重数”问题解决了，下一步才按次数分组。

\subsection{第二步：按次数分组（distinct-degree）}


从这一步起假设 $f$ 无平方因子。

\textbf{引理（关键）。} 设 $f$ 无平方因子，则

\[
\gcd\bigl(f,\ x^{q^i}-x\bigr)
\]

等于“$f$ 的那些次数整除 $i$ 的不可约因子”之积。

\textbf{证明。} 由上一节，$x^{q^i}-x=\prod_{e\mid i}\Phi_e$，而 $\Phi_e$ 是全体 $e$ 次首一不可约多项式之积。所以“$f$ 的某个不可约因子 $r$ 整除 $x^{q^i}-x$”当且仅当 $\deg r\mid i$。又 $f$ 无平方因子，故 $\gcd$ 就是这些因子的乘积。$\square$

\textbf{算法（渐进除法版）。}

\begin{itemize}
\item 初始化：令 $f^\ast\leftarrow f$，$i\leftarrow1$。
\item 循环（当 $\deg f^\ast\ge 2i$ 时）：计算 $g=\gcd\bigl(f^\ast,\ x^{q^i}-x\bigr)$。若 $g\ne1$，就输出 $g$——它是 $f^\ast$ 中全体 $i$ 次不可约因子之积（次数相同，交给第三步）；并令 $f^\ast\leftarrow f^\ast/g$。然后 $i\leftarrow i+1$。
\item 收尾：循环结束后，若 $\deg f^\ast>0$，输出 $f^\ast$，它本身不可约。
\end{itemize}


\textbf{正确性。} 用归纳不变量：每轮开始时，$f^\ast$ 无平方因子，且它的所有不可约因子的次数都 $\ge i$。

\begin{itemize}
\item $i=1$ 时显然（次数都 $\ge1$）。
\item 设第 $i$ 轮开始时不变量成立。由引理，$g=\gcd(f^\ast,x^{q^i}-x)$ 是“$f^\ast$ 中次数整除 $i$ 的因子”之积；而不变量说次数都 $\ge i$，所以“次数整除 $i$”等价于“次数等于 $i$”。于是 $g$ 恰是全体 $i$ 次因子之积；除掉它后，剩余因子的次数都 $>i$，即 $\ge i+1$，不变量保持。
\item 退出时 $\deg f^\ast<2i$，而 $f^\ast$ 的每个因子次数都 $\ge i$。若 $f^\ast$ 有两个或更多因子，则 $\deg f^\ast\ge2i$，矛盾。所以 $f^\ast$ 至多一个因子，即不可约（或为 $1$）。$\square$
\end{itemize}


\textbf{为什么不能用累积的 $\gcd$ 直接相除？} 一个常见的错误写法是：令 $G_i=\gcd(f,x^{q^i}-x)$，然后认为 $G_i/G_{i-1}$ 是次数恰为 $i$ 的因子。\textbf{这不成立}，因为 $G_i$ 不嵌套。反例（$\mathbb F_3$）：

\[
f=(x+1)(x^2+1)(x^3-x+1)=x^6+x^5+x^3+1 .
\]

此时

\[
G_1=x+1,\qquad G_2=(x+1)(x^2+1)=x^3+x^2+x+1,
\]
\[
G_3=(x+1)(x^3-x+1)=x^4+x^3-x^2+1 .
\]

$G_2$ 不整除 $G_3$（$G_3$ 含 $x^3-x+1$，$G_2$ 含 $x^2+1$，两者互不相同），所以 $G_3/G_2$ 没有意义。原因：$G_2$ 收集“次数整除 $2$”的因子，$G_3$ 收集“次数整除 $3$”的因子，两个集合互不包含。\textbf{正确做法是每一步先把已找到的因子除掉}（渐进版），使剩余因子的次数都 $\ge i$。

\textbf{例三（$\mathbb F_3$ 上的完整分组）。} 仍取上面的 $f=x^6+x^5+x^3+1$。

\begin{itemize}
\item $i=1$：$f^\ast=f$，$\deg f^\ast=6\ge2$。$g=\gcd(f,x^3-x)=x+1$。输出 $x+1$；$f^\ast=(x^2+1)(x^3-x+1)$，次数 $5$。
\item $i=2$：$\deg f^\ast=5\ge4$。$g=\gcd(f^\ast,x^9-x)=x^2+1$。输出 $x^2+1$；$f^\ast=x^3-x+1$，次数 $3$。
\item $i=3$：$\deg f^\ast=3<6=2i$，退出；输出 $f^\ast=x^3-x+1$，它不可约。
\end{itemize}


最终得到三个因子：$x+1$、$x^2+1$、$x^3-x+1$。

\textbf{例四（$\mathbb F_2$）。} 取 $f=x^3+1=(x+1)(x^2+x+1)$。$i=1$：$g=\gcd(f,x^2-x)=x+1$，输出后 $f^\ast=x^2+x+1$（次数 $2$）；$i=2$ 时 $2\ge4$ 不成立，退出，输出 $x^2+x+1$。其中 $x^2+x+1$ 不可约。

\textbf{计算 $x^{q^i}\bmod f^\ast$。} 直接展开 $x^{q^i}$ 是不可能的（次数指数级）。用快速幂：先在 $\mathbb F_q[x]/(f^\ast)$ 中算 $x^q$，再反复“取 $q$ 次幂”：

\[
x^{q^i}=\bigl(x^{q^{i-1}}\bigr)^{q}.
\]

每次取 $q$ 次幂可以做 $O(\log q)$ 次多项式乘法再取模，因此第二步的总代价是 $\deg f$ 的多项式级别。

\subsection{第三步：同次数拆分（Cantor–Zassenhaus）}


现在假设 $f$ 无平方因子，而且它的全部不可约因子次数都相同，记这个公共次数为 $d$：

\[
f=r_1r_2\cdots r_k,\qquad \deg r_1=\cdots=\deg r_k=d .
\]

\textbf{中国剩余定理拆环。} 不同的不可约多项式两两互素，所以由中国剩余定理（Lec09 第 3 节），

\[
\mathbb F_q[x]/(f)\ \cong\ \mathbb F_q[x]/(r_1)\times\mathbb F_q[x]/(r_2)\times\cdots\times\mathbb F_q[x]/(r_k).
\]

每个 $\mathbb F_q[x]/(r_i)$ 都是域；因为 $\deg r_i=d$，它同构于

\[
\mathbb F_{q^d}.
\]

于是“在 $\mathbb F_q[x]/(f)$ 里算东西”，等价于“在每个 $\mathbb F_{q^d}$ 里分别算，再拼起来”。

\textbf{关键事实（$q$ 为奇数时）。} 在 $\mathbb F_{q^d}$ 中取非零元 $y$。由 Fermat 小定理的有限域版，

\[
y^{q^d-1}=1 .
\]

因为 $q$ 是奇数，$q^d-1$ 是偶数，令 $m=(q^d-1)/2$。则

\[
\bigl(y^{m}\bigr)^2=y^{q^d-1}=1,
\]

而在域中 $z^2=1$ 只有 $z=\pm1$，所以

\[
y^{(q^d-1)/2}=\pm1 .
\]

更进一步，\textbf{$\mathbb F_{q^d}^\times$ 中恰好一半元素给出 $+1$，一半给出 $-1$}：方程 $z^{(q^d-1)/2}-1=0$ 在 $\mathbb F_{q^d}$ 中至多 $(q^d-1)/2$ 个根，$z^{(q^d-1)/2}+1=0$ 也至多 $(q^d-1)/2$ 个；而 $q^d-1$ 个元素每个都满足其中一个，所以两边各恰好占一半。

\textbf{随机采样与分裂。} 均匀随机取一个多项式 $s\in\mathbb F_q[x]/(f)$（即随机取次数 $<\deg f$ 的系数）。由 CRT，它对应一个随机元组

\[
s\ \longleftrightarrow\ (s_1,\dots,s_k),\qquad s_i=s\bmod r_i\in\mathbb F_{q^d},
\]

而且因为 CRT 是同构（两边元素个数相同），均匀随机的 $s$ 对应\textbf{独立均匀}的 $s_i$。

\begin{itemize}
\item 若某个 $s_i=0$：那么 $r_i\mid s$，直接算 $\gcd(s,f)$ 就可能切出因子（这一步顺手处理）。
\item 否则每个 $s_i\ne0$，令 $\varepsilon_i=s_i^{(q^d-1)/2}\in\{\pm1\}$。计算
\end{itemize}


\[
h=s^{(q^d-1)/2}-1\pmod f .
\]

在分量 $i$ 上，$h\bmod r_i=\varepsilon_i-1$：当 $\varepsilon_i=1$ 时 $r_i\mid h$，当 $\varepsilon_i=-1$ 时 $r_i\nmid h$。因此

\[
\gcd(h,f)=\prod_{\varepsilon_i=1}r_i ,
\]

只要 $\varepsilon_i$ 不全相同，这个 $\gcd$ 就是 $f$ 的\textbf{非平凡因子}。

\textbf{成功率。} 需要先处理两个“退化情况”：

\begin{itemize}
\item 若 $s=0$（所有分量都为 $0$），则 $\gcd(s,f)=f$，得不到信息，重抽即可；这个概率是 $1/q^{\deg f}$，可以忽略。
\item 若某个 $s_i=0$ 但 $s\ne0$，则 $r_i\mid s$，直接用 $\gcd(s,f)$ 就能切出因子（不必再算 $h$）。
\end{itemize}


排除这些情况后，每个 $\varepsilon_i$ 独立地以 $1/2$ 的概率取 $+1$、以 $1/2$ 的概率取 $-1$。试验失败当且仅当全部 $\varepsilon_i$ 相同，概率为

\[
\left(\frac12\right)^k+\left(\frac12\right)^k=2^{1-k}\le\frac12\qquad(k\ge2),
\]

所以一次试验成功（得到非平凡因子）的概率至少 $1/2$。失败只有两种情形：全部 $+1$（此时 $\gcd=f$）或全部 $-1$（此时 $\gcd=1$）。

\textbf{严格合并一步。} 设 $F$ 为失败事件。失败只可能来自两类抽样：$s=0$（此时 $\gcd(s,f)=f$、$h=-1$，什么也得不到），或所有分量都非零且符号全同。因此

\[
F\subseteq\{\text{所有分量非零且符号全同}\}\cup\{s=0\},
\]

\[
P(F)\le\Bigl(1-\frac1{q^d}\Bigr)^k 2^{1-k}+\frac1{q^{kd}} .
\]

$k\ge2$ 时这个上界严格小于 $1/2$（$k=2$ 时它等于 $1/2-1/q^d+3/(2q^{2d})<1/2$；$k\ge3$ 时更小），故 $P(\text{成功})\ge1/2$。这里把「某个分量为 $0$ 而 $s\ne0$」的情形算作成功（用 $\gcd(s,f)$ 直接拆分）；若完全不处理退化、只算 $\gcd(s^{(q^d-1)/2}-1,f)$，$k=2$ 时的成功率是 $1/2-1/(2q^{2d})$（本讲例子 $40/81$）。

\textbf{算法。}

\begin{itemize}
\item 若 $\deg f=d$，则 $f$ 不可约，返回。
\item 随机取 $s$。先算 $g_0=\gcd(s,f)$：若 $1<\deg g_0<\deg f$，说明有某个分量 $s_i=0$，直接拆分 $f=g_0\cdot(f/g_0)$。
\item 否则计算 $t=s^{(q^d-1)/2}-1\bmod f$，再算 $g=\gcd(t,f)$。
\item 若 $1<\deg g<\deg f$：递归分解 $g$ 与 $f/g$（两者仍然无平方因子、因子次数都是 $d$）。
\item 否则重复上一步。
\end{itemize}


期望试验次数是常数（约 $2$ 次）；每次分裂至少把因子个数分开一些，总期望试验次数与 $k$ 同阶。

\textbf{完整手算例子（$\mathbb F_3$）。} 取

\[
f=(x^2+1)(x^2+x-1)=x^4+x^3+x-1,\qquad q=3,\ d=2,
\]

于是 $(q^d-1)/2=(9-1)/2=4$。

\textbf{成功的一次：取 $s=x$。} 先算 $x^4\bmod f$：由 $f=0$ 得 $x^4=-x^3-x+1$。所以 $s^4\bmod f=-x^3-x+1$，而

\[
h=s^4-1\equiv -x^3-x=-(x^3+x)\pmod f .
\]

于是 $\gcd(f,h)=\gcd(x^4+x^3+x-1,\ x^3+x)$。直接观察：$x^3+x=x(x^2+1)$，而 $x^2+1\mid f$，故

\[
\gcd(f,h)=x^2+1 .
\]

一次就拆开：$f=(x^2+1)(x^2+x-1)$。

\textbf{失败的样子。} 取 $s=x+1$：计算得 $s^4\equiv-1\pmod f$，于是 $h\equiv-2\equiv1$（常数），$\gcd=1$，失败。取 $s=x^2$：计算得 $s^4\equiv1\pmod f$，于是 $h\equiv0$，$\gcd=f$，也失败。其余随机选择中，只要两个分量的 $\pm1$ 不同，就能拆开；因此成功率超过 $1/2$。

\textbf{$q$ 为偶数时怎么办？} 若 $q=p^r$ 且 $p=2$，则 $q^d-1$ 是奇数，$(q^d-1)/2$ 不是整数，上面的“正负号拆分”失效。替代做法是用\textbf{迹映射}

\[
T(y)=y+y^q+y^{q^2}+\cdots+y^{q^{d-1}}\in\mathbb F_q
\]

把 $\mathbb F_{q^d}^\times$ 按 $T$ 的值分组，再对 $\gcd(T(s)-a,f)$ 做拆分。原理与奇数情形相同（每个分量独立、以一定概率落在不同值上），细节不要求。手写稿注明的“暂且只考虑 odd $q$”正是这个原因。

\subsection{三个算法合起来}


\begin{itemize}
\item \textbf{去平方因子}：解决“重数”问题，输出若干个无平方因子部分；
\item \textbf{按次数分组}：把次数不同的不可约因子分开，输出按次数分好的乘积（每堆里次数相同）；
\item \textbf{同次数拆分}：把同次数、互不相同的因子逐个分开，输出单个不可约因子。
\end{itemize}


对每个无平方因子部分依次做后两步；每一步切出的因子再递归调用整条流水线。

\textbf{复杂度结论（不必背具体阶）。} 三个步骤都只用到多项式乘法、快速幂、辗转相除和中国剩余定理，在 $\mathbb F_q[x]$ 中都是多项式时间；这与整数分解“目前没有已知多项式时间算法”形成对比。课堂最后顺带提到量子算法（Shor 算法）能在多项式时间分解整数，但那是另一条线，本课程不展开。

\subsection{本节的勘误与提醒（对应手写稿第 3 页下半页）}


\begin{itemize}
\item \textbf{手写稿第一块（square-free）的例子是 $x^{2p}-x^p+3=(x^2-x+3)^p$。} 注意例子里的“$3$”要在 $\mathbb F_p$ 中理解（$3^p=3$），并且 $p=2$ 时也成立（此时 $-x^p=x^p$，$3\equiv1$）。写 $x^2+x+3$ 还是 $x^2-x+3$ 只差一个符号，在特征 $2$ 下相同；一般 $p$ 下应为 $x^2-x+3$。
\item \textbf{手写稿第二块（distinct-degree）那两行只是“单步说明”。} $r_1\mid x^{q^{d_1}}-x$、$r_2\nmid x^{q^{d_1}}-x$（其中 $d_1=\deg r_1$）说明不同次数的因子可以用 $\gcd$ 分开；完整算法必须按上文的渐进除法写，不能直接用 $G_i/G_{i-1}$。
\item \textbf{“some degree” 是“同次数”（equal-degree）的意思}，不是“某个次数”。手写稿标题里的 some degree 即 Cantor–Zassenhaus 这一步。
\item \textbf{“暂且只考虑 odd $q$”是必要的限制。} $q$ 偶数时 $(q^d-1)/2$ 不是整数，正负号拆分不成立，需要改用迹映射。
\item \textbf{公式里的 $d$ 是“每个不可约因子的公共次数”，不是 $\deg f$。} 例如 $f$ 是 6 个二次因子之积时，$d=2$，指数用 $(q^2-1)/2$。
\item \textbf{$\gcd$ 全程在 $\mathbb F_q[x]$ 里做。} 辗转相除不会把系数带到扩域；这也是“用 $x^{q^i}\bmod f^\ast$ 去算 gcd”合法的原因。
\item \textbf{$f'=0$ 不是“$f$ 没有重根”，恰恰相反：它是 $p$ 次幂，重数全都很高。} 判断重因子要用 $\gcd(f,f')$，不是看 $f'$ 是否为零。
\item \textbf{平方因子与不可约是两件事。} 去平方因子后的 $f$ 可能仍然可约（如 $x^4+x^3+x-1=(x^2+1)(x^2+x-1)$ 无平方因子但可约）。
\item 课堂转写里“精确讲，这个地方应该是……”一处老师在纠正口误，指的就是 $f'$ 与 $\gcd$ 的关系；以本节引理为准。
\end{itemize}


\section*{一页式总览}
\begin{itemize}
\item $K/F$ 是扩域，$[K:F]=\dim_F K$ 是次数；$K$ 是 $F$-线性空间，这是整讲的计算引擎。
\item 造域的三种方式：$F[x]/(f)$（添形式根）、$\operatorname{Frac}(F[x])=F(x)$（添分母）、$F(\alpha)$（在具体扩域里取真实根）。
\item $f$ 不可约、$f(\alpha)=0$ 时 $F[x]/(f)\cong F(\alpha)$，$\bar x\mapsto\alpha$；$[F(\alpha):F]=\deg m_{\alpha,F}$。
\item 域同态非单即零；有限扩张必代数，反之不然；超越元满足 $\operatorname{Frac}(F[x])\cong F(\alpha)$。
\item 塔律 $[H:F]=[H:K][K:F]$；中间域次数整除总次数。尺规作图每步至多二次扩张，可作数满足 $[\mathbb Q(x):\mathbb Q]=2^m$。
\item 倍立方需 $\sqrt[3]{2}$（次数 $3$），三等分 $60^\circ$ 需 $\cos20^\circ$（次数 $3$），都推出 $3\mid2^m$ 的矛盾；化圆为方：若可作则 $\sqrt\pi$ 必在 $\mathbb Q$ 上代数；但 $\pi$ 超越 $\Rightarrow$ $\sqrt\pi$ 超越，矛盾。
\item 分裂域：$f$ 在 $E$ 上完全分裂且 $E=F(\text{全部根})$；存在，且 $[E:F]\le(\deg f)!$；在同构意义下唯一。
\item 有限域：$\mathbb F_{p^d}$ 存在唯一；$\mathbb F_{p^d}$ 是 $x^{p^d}-x$ 在 $\mathbb F_p$ 上的分裂域；$x^{q^d}-x=\prod_{e\mid d}\Phi_e$，$\deg\Phi_d=dN_d$，$q^d=\sum_{e\mid d}eN_e$。
\item 每个次数 $d$ 都有不可约多项式：$\deg\Phi_d\ge q^d-\sum_{e\mid d,e<d}q^e>0$；$N_d=\frac1d\sum_{e\mid d}\mu(d/e)q^e$。
\item 有限域上因式分解：去平方因子（$\gcd(f,f')$；$f'=0$ 时开 $p$ 次方）$\to$ 按次数分组（渐进除法，$f^\ast$ 与 $x^{q^i}-x$ 做 $\gcd$）$\to$ 同次数拆分（Cantor--Zassenhaus：随机 $s$，算 $\gcd(s^{(q^d-1)/2}-1,f)$；$q$ 奇且 $k\ge2$ 时，按含退化处理的算法，成功概率 $\ge1/2$）。
\end{itemize}

\section*{参考资料}
\begin{itemize}
\item \textbf{[DF]} D. S. Dummit, R. M. Foote. \emph{Abstract Algebra}, 3rd ed. §13.1--13.4（域扩张、代数扩张、古典尺规作图、分裂域），§14.3（有限域与不可约多项式计数、Möbius 公式）。有限域在第 14 章 §14.3，不在 §13.4。
\item \textbf{[DS]} 邓小铁、王畅。《离散数学与结构》§10.3（特征、素域、$F[x]/(f)$ 是域 $\iff f$ 不可约、有限域的存在唯一性）。该书没有分裂域与多项式分解算法的系统讨论。
\item \textbf{[讲义]} 课程讲义 \texttt{25-26-lecture\_notes.pdf} 第八章：§8.1 域扩张（定义 8.1--8.2、定理 8.1--8.3、命题 8.1--8.4）、§8.2 分裂域（定义 8.3--8.4、定理 8.4）、§8.3 有限域（定理 8.5、命题 8.5--8.6、引理 8.1、定理 8.6--8.9）。该讲义为 2025--2026 秋季版本。
\item 原始材料：本地归档（两段转写与三页手写原稿），不纳入课程仓库。
\item 第 3 页的零基础讲解见 \texttt{专题整理/图离散Lec11-笔记第3页讲解稿.pdf}。
\end{itemize}
\end{document}
